% Zone „Kennst du schon" – Fokus Nullstellen (quadratische-funktionen, Einheit 4)
\einheitenkopf[zone][Kennst du schon]{Das kennst du schon}
\setcounter{aufgabe}{0}

\begin{aufgabe}{Ich kann die Nullstelle einer Geraden berechnen.}
Berechne die Nullstelle der Geraden.
\nullstellenliste{2x-8, -3x+6}
\end{aufgabe}

\begin{aufgabe}{Ich kann Zahlen quadrieren, auch negative Zahlen.}
Berechne.
\begin{teilezwei}
\tz $(-6)^2 = $ \leerfeld & \tz $(-1{,}5)^2 = $ \leerfeld \\
\end{teilezwei}
\end{aufgabe}

\begin{aufgabe}{Ich kann eine Quadratwurzel ziehen und runden.}
Berechne.
\begin{teile}
\teil $\sqrt{81} = $ \leerfeld
\anweisung{Runde auf zwei Stellen nach dem Komma.}
\teil $\sqrt{13} \approx $ \leerfeld
\end{teile}
\end{aufgabe}

\begin{aufgabe}{Ich kann eine Gleichung der Form $x^2 = $ Zahl lösen.}
Löse die Gleichung. Gib alle Lösungen an.
\begin{gleichungsraster}[2]
\gl{x^2 = 64} & \gl{x^2 = -9} \\
\end{gleichungsraster}
\end{aufgabe}

\begin{aufgabe}{Ich kann eine quadratische Gleichung mit der p-q-Formel lösen.}
Löse die Gleichung.
\begin{gleichungsraster}[4]
\gl{x^2 + 8x + 15 = 0} & \gl{x^2 - 2x - 24 = 0} \\
\end{gleichungsraster}
\end{aufgabe}
